% SPDX-License-Identifier: GPL-3.0-or-later
\documentclass[11pt]{article}
\usepackage[T1]{fontenc}
\usepackage{lmodern,amsmath,amssymb,amsthm,booktabs,hyperref}
\newtheorem{theorem}{Theorem}[section]
\newcommand{\result}{\par\noindent\textit{Result.}\ }
\usepackage[margin=1in]{geometry}
\title{PasTeX User's Guide}
\author{PasTeX contributors}
\date{Version 0.10.0-rc.2 (release candidate)}
\begin{document}
\maketitle
\tableofcontents

\section{Your first editing session}
PasTeX edits UTF-8 \LaTeX{} files directly. Save a copy of this guide to
practice the editing steps and the source/result examples in the appendix.
Installation, build commands and the example index are in \texttt{docs/developer-guide.tex}.

\begin{enumerate}
\item Open a file with Ctrl+O. Click a word to place the caret in the instant editor.
\item Type a sentence. Press Enter to start another paragraph; an unused empty
paragraph disappears when you leave it.
\item Drag to select characters, double-click a word, or use Shift with arrow keys.
Use Ctrl+C, Ctrl+X and Ctrl+V to copy, cut and paste. Selections can cross paragraphs.
\item Save with Ctrl+S. Undo with Ctrl+Z and redo with Ctrl+Y. The narrow arrows
beside those buttons show histories; selecting an entry applies several steps.
\item Open another file. Drag its tab to reorder it, or use the tab's right-click
menu. Close a tab with its close button or Ctrl+W; unsaved changes prompt for saving.
\end{enumerate}

\subsection{Find a command}
Home contains text, paragraph and editing tools. Insert adds objects; References
handles labels and bibliography. Contextual tabs appear for the current equation,
table, figure, theorem, box, listing or command definition. View controls the
workspace; Settings controls the document and application.
Hover over an icon for its description. Unavailable buttons are gray.
Use Alt+Q to search commands by name. Ctrl+F1 collapses or restores the ribbon.
The traditional menubar and context menus provide the same actions in task groups.
Home, Editing, Properties opens the settings for the command or object at the
caret. Its argument editor supports required, optional and unbraced arguments,
with source highlighting and controls to add, remove or reorder them.

\subsection{Focus and paper pages}
Press F11, or choose View, Focus/fullscreen, to hide the ribbon, tabs and
navigation. Escape or Exit focus restores the previous window and controls.
The status-bar layout buttons switch between Continuous and Paged layout in the
main window; View, Paged layout (Ctrl+Alt+P) is the same command. Click a word on
native paper pages to edit text and formulas instantly. The shared renderer
uses the document paper size and margins, breaks long paragraphs across pages
and reflows changed content through the native renderer. Save and the
formatting ribbon work in the main window. Ctrl+R opens the compiled PDF for
checking final pagination, floats and package-specific output routines.
Native pages support common single-column float placement and page styles,
including page numbers, fancyhdr fields, marks and header/footer rules.
Selected text stays synchronized with Full document source across paragraphs.
The main status bar shows the current page and total.
Sources start collapsed in page view. Hover to expand them; temporary viewing
space holds source, labels, comments and active command/environment names
while preserving page assignments. Collapsing source or leaving a command
removes its extra space.
Visual editor and Source editor switch the active paragraph within either layout.
Continuous layout restores the paragraph workspace.

\subsection{Source and compiled PDF side by side}
Choose View, Source / compiled PDF (Ctrl+Shift+P), or its status-bar button.
The complete editable source is on the left and the compiled PDF is on the
right. Drag the divider to resize the panes. Source fonts, highlighting and
completion use the normal editor settings; Ctrl+S saves the document.

Compilation runs after a short typing pause. Refresh PDF, above the PDF pane,
refreshes explicitly; errors leave the previous PDF visible. The right arrow
on the divider locates the source caret in the PDF. The left arrow locates the
visible PDF position in source. Ctrl+click the PDF to locate a specific line.
These actions require the SyncTeX utility supplied by your \TeX{} distribution.
Open the master file to compile a project; included files are read from disk.
Use Continuous layout or Paged layout to return to native visual editing.

\subsection{Beamer presentations}
Open \texttt{examples/beamer.tex} to try native slide editing. The document class
selects slide dimensions and text defaults. Continuous layout exposes overlay
contents for editing; Paged layout repeats each frame for its overlays. Edits
change the shared source, including other overlays of that frame. Fit Width and
Fit Page use the slide aspect ratio.

For example, \verb|\only<2>{Second slide}| occupies space on overlay two only,
whereas \verb|\uncover<2->{Revealed text}| reserves space on earlier overlays.
\verb|\pause| advances the reveal step; \verb|\alert<2>{Important}| changes the
text color on overlay two. The native renderer also handles \verb|\onslide|,
\verb|\visible|, \verb|\invisible|, \verb|\alt| and \verb|\temporal|, plus ordinary
columns, blocks and theorem environments. An \texttt{overprint} environment
replaces alternatives in the same area. Theme font/color declarations and their
options are read from installed theme files and document overrides. Infolines
footers show document metadata and frame counts.

Use \verb|\titlepage| for document metadata and \verb|\framezoom<1><2>(0cm,0cm)(5cm,3cm)|
for an enlarged region on the second overlay. Handout mode combines the ordinary
overlays. The example includes both features; command suggestions also complete
frames, block titles and overlay arguments.

Beamer support is experimental. Use Ctrl+R to verify theme templates, notes,
zoom hyperlinks and automatic frame breaks in the compiled presentation.

\subsection{Import and exchange documents}
Use Insert, Spreadsheet data, to read XLSX, ODS, Gnumeric, CSV or TSV. Choose
a sheet and an optional range such as \texttt{B2:D20}; inspect the grid, then
insert a table. Import reads cached spreadsheet values into a LaTeX table.

File, Import document, converts Word, ODT, RTF or LyX to a new \texttt{.tex} file.
Pandoc preserves richer DOCX content and extracts media when installed; the
built-in DOCX reader imports text, headings, bold/italic runs and simple tables.
Legacy Word, ODT and RTF require LibreOffice; LyX requires its own executable.
The original file stays intact. Check converted content before publishing.

\subsection{Tools and shortcuts}
File, Export project ZIP, packages the source and assets in the document folder.

Settings, Options, \TeX{}, configures each engine's executable and arguments.
Use \texttt{\{source\}} for the generated source and \texttt{\{output\}} for
its temporary directory, then Save compile command. Restore defaults resets
the fields. Tools, Keyboard shortcuts, assigns or clears command bindings;
clear a conflicting binding first. Escape always leaves focus mode.

\subsection{Notes}
Insert, Notes, adds a footnote, a separate footnote mark/text pair or a margin
note. Edit note text and mathematics directly in instant preview. The continuous
view places editable notes below their paragraph; compiled pages use \TeX{}'s
footnote and margin placement. See \texttt{text-formatting.tex}.

\subsection{Understand the previews}
Settings, Preview, \emph{Compile continuous preview automatically} controls
background \LaTeX{} compilation. Turn it off to keep continuous view native;
Refresh preview compiles once on request. Native text and mathematics use bundled
fonts and metrics for standalone editing. PDF preview,
PDF images and export need the corresponding optional tools. Unsupported package
programs remain editable source boxes. Cached drawing preview falls back to the
native drawing backend when the compiler is unavailable.
Use Settings, System, Check environment and dependencies (also in Help) to
inspect detected runtimes, compiler paths and optional tools. Select a missing
component for installation instructions; Check again refreshes the report.

Instant preview updates locally while you type. Permanent preview uses your
selected \TeX{} engine. During compilation, changed content has a yellow frame;
a red frame means compilation failed. Other paragraphs retain their last preview.
The instant editor remains usable after an error. Inspect Settings, Compiler log,
correct the source, and refresh. Ctrl+R opens the compiled full-document PDF;
File, Export PDF saves it to a chosen path.

The continuous editor uses a paragraph layout. The PDF is the
reference for pagination, floats and arbitrary package programs. Native package
support is partial; unsupported commands and environments retain editable source
boxes. Diagram preview can use cached \TeX{} output when native drawing differs.

\section{Text, headings and lists}
\begin{enumerate}
\item Select a phrase and choose Home, Font, Bold or Italic. Click again to remove
that style. The buttons and family/size lists reflect the selection.
\item Choose Roman, Sans Serif or Monospace to change its family. Small caps,
slanted and upright styles can be combined. Clear formatting removes character styles.
\item Use Text color for a color; its arrow restores the document default.
Named sizes follow the document class. Custom accepts a number and unit.
\item Place the caret in a paragraph and choose its heading in Home, Paragraph.
Normal removes the heading wrapper. Available headings follow the document class;
declared theorem environments also appear in this list.
The \texttt{secnumdepth} counter controls numbering; value 5 includes
\verb|\paragraph| and \verb|\subparagraph|, as shown in \texttt{text-formatting.tex}.
\item Choose a list from Insert. Enter adds an item, and the context menu edits
list structure. Only one unused empty item is retained.
\end{enumerate}

Try nested styles, alignment, literal code and escaped symbols in
\texttt{text-formatting.tex}. In visual text, percent and other reserved symbols
are escaped when inserted. Type \verb|~| for a nonbreaking space, as in
\verb|Equation~\eqref{eq:guide}|. Ctrl+/ comments or uncomments selected text;
with an empty selection it affects the current source line. Comments have green
editing boxes. Math braces and explicit line/space commands have editing marks
used for editing. Click a horizontal or vertical space marker to
choose a preset or enter its size.

\section{Write and edit mathematics}
\begin{enumerate}
\item Press Alt+= in body text to start an inline equation, or choose Inline or
Displayed from Insert. Selecting a whole equation and pressing Alt+= removes
its delimiters. The Equation tab offers inline/display conversion too.
\item Insert a fraction, root or matrix from Equation. Click an empty argument's
placeholder and type. Placeholders guide editing and disappear as content is entered.
\item Use Tab or Shift+Tab between fields. Arrow keys move through the expression.
Up/Down moves between a fraction's numerator and denominator, retaining the
horizontal caret position, before moving to another equation row.
Space leaves the current argument or script; at the formula's end it returns to
body text. Inside \verb|\text{...}|, Space inserts ordinary text spaces.
\item Backspace at the start of a script or brace group removes its wrapper and
keeps the contents. Backspace in an empty equation removes the equation.
Immediately after a command, Backspace selects the command; press it again to
delete the selection. Inside its argument, Backspace deletes text normally.
Empty superscripts and subscripts disappear when the caret leaves them.
\item Use the symbol galleries for Greek letters, arrows, operators, functions,
accents and math alphabets. Labelled arrows have editable argument fields.
\item For aligned equations, choose an environment from Insert. Enter and the
context menu add/remove rows. The Equation tab and context menu duplicate the
current row or column; column actions are disabled for fixed-column environments
such as gather and cases. Matrix tools also change delimiters.
\end{enumerate}

For example, enter this in equation source:
\begin{verbatim}
\sum_{n=1}^{\infty}\frac{1}{n^2}=\frac{\pi^2}{6}
\end{verbatim}
\begin{equation}\label{eq:guide}
\sum_{n=1}^{\infty}\frac{1}{n^2}=\frac{\pi^2}{6}.
\end{equation}
Use \texttt{math.tex} to practice structures, fonts, accents and environments.
Copying mathematics into body text supplies math delimiters. Copying words from
\verb|\text{...}| omits its wrapper. Typing a backslash opens command hints:
Up/Down chooses a hint, Enter or Tab accepts it, and Escape cancels. Space accepts
the typed command and leaves its argument ready for editing.
Suggestions follow the text/math context and definitions before the caret.

\subsection{Shortcuts for LyX users}
The instant editor accepts common sequences from
\href{https://git.lyx.org/repos/lyx.git}{LyX's keyboard bindings}.
Press and release Alt+M, then type the keys in the second column.
Escape cancels a sequence; the status bar shows the pending choices.
These sequences also work in page view and leave source editors unchanged.
\begin{center}
\begin{tabular}{ll}
Action&Keys after Alt+M\\\hline
Fraction / square root / indexed root&\texttt{f / s / r}\\
Superscript / subscript&\texttt{e / x}\\
Text inside a formula&\texttt{m}\\
Sum / integral / contour integral / infinity&\texttt{u / i / y / 8}\\
Paired delimiters&\verb|(|, \verb|[|, \verb|{|, \verb|<|, \verb.|.\\
Hat / vector / overline / dot / double dot&\verb|h / v / b / . / "|\\
Greek letters&\texttt{g a} ($\alpha$), \texttt{g p} ($\pi$), \texttt{g G} ($\Gamma$)\\
Insert / delete / duplicate row&\texttt{w i / w d / w c}\\
Insert / delete / duplicate column&\texttt{c i / c d / c c}\\
Inline / numbered equation / align&\texttt{t i / t d / t a}\\
\end{tabular}
\end{center}
Fractions, roots, accents and delimiters enclose selected content. Column actions
require a suitable matrix or alignment. Ctrl+M starts inline math, or inserts
\verb|\text{...}| within mathematics. Ctrl+Shift+M creates or converts to displayed math.
Ctrl+Shift+Z is an additional Redo shortcut.

\section{Move between visual editing and source}
Ctrl+L opens selected source in a gray box in place. Hover over a body command,
equation, drawing or group such as \verb|{\itshape Emphasized words}| to reveal
a dashed outline and plus button. Wrapped outlines follow the lines, with the
button to the left of the second line. It expands source below the preview; the minus
on the dashed divider collapses it, including inside \verb|\fbox|. Body groups
and equations start collapsed in both layouts;
figures with source editors start expanded in continuous layout. Other content stays visible.
Hover in the left gutter of a list or other environment to expand its complete source.
The source editor uses the font, wrapping and highlighting from Settings.
Drag its bottom grip to resize it; Fit content restores automatic sizing.
Source editors suggest commands as you type; Enter or Tab accepts an item,
Escape dismisses the list, and Ctrl+Space opens it again. For example, accept
\verb|\begin{}|, type \texttt{fig}, then accept \texttt{figure}. PasTeX adds
the matching \verb|\end{figure}| and places the caret on an indented line inside.
Suggestions include loaded package commands and definitions preceding the caret.
Comments and verbatim text remain literal.

After visual edits, opening source or leaving the paragraph formats supported
environment indentation and equation rows. This preserves comments, literal
text and the caret position. Plot expressions and data-set function definitions
use the same source editor, including wrapping and syntax highlighting.

Ctrl+Shift+U opens the complete source in a separate, modeless form. Preamble,
body and ending are editable there. Caret and selection synchronize with visual
editing. After scrolling independently, use these buttons:
\begin{description}
\item[Go to preview position] Scroll full source to the preview's visible caret,
or to the visible paragraph when the caret is off screen.
\item[Show in preview] Scroll the document to the visible source caret, or to the
top visible source line when the caret is off screen.
\end{description}
The same actions are in View, Locate. They preserve the text and selection while moving the view.

\section{Configure a document and define commands}
Open Settings, Preamble and document settings. Choose the document class and
edit its options as flags or key/value rows. Add, remove and reorder packages;
each has its own options and related setup commands. Document information sets
title, author and date. Complete preamble exposes the remaining source with the
same theme as other editors. Apply updates the document.

The Fonts tab configures Latin and CJK families, named font commands and font
features. Choose an installed font or an OpenType/TrueType file; use options such
as \texttt{BoldFont}, \texttt{ItalicFont}, \texttt{BoldItalicFont}, \texttt{Path},
\texttt{Extension} and \texttt{Scale}. CJK font commands require Xe\LaTeX;
Latin fontspec commands also work with Lua\LaTeX. See \texttt{chinese.tex}.
Feature defaults can target font names or family commands such as
\verb|\rmfamily|; Append retains the preceding defaults. Named families and
single-face commands can also be configured here.
File paths are relative to the document directory. Native preview uses the
specified file's outlines, including explicit bold and italic faces and CJK fonts.
Engine tests such as \verb|\ifxetex|, \verb|\ifluatex| and \verb|\ifpdftex|
select their active branch while preserving the source; nested \verb|\else|
and \verb|\fi| are supported. \verb|\protect| preserves normal typesetting.
For stretchable spacing, try \verb|Left\hfill Right| or
\verb|Contents\dotfill Page 1| in \texttt{text-formatting.tex}.
Native preview uses the selected font's outlines and HarfBuzz shaping for
kerning, ligatures and common OpenType features, including small capitals and
old-style numbers. It also applies numeric scale and word-space settings.
Other feature options remain in the source for compilation; complex-script,
variable-font and vertical-layout typography use the compiled \TeX{} output.

\texttt{text-formatting.tex} demonstrates editable \texttt{minipage} columns,
\verb|\hrule|, \verb|\rule| and \verb|\hdashrule| (from \texttt{dashrule}).
Minipages accept width, outer alignment, height and inner alignment options.
Click a rule to switch its command or change its dimensions and dash pattern;
choosing a dashed rule adds the required package.
Margin notes remain editable outside the text column. Named font switches,
such as \verb|\examplekai| in \texttt{chinese.tex}, have gray source markers.

In Packages, select a package and choose Edit source (.sty). PasTeX searches the
document folder first, then the installed \TeX{} distribution. The source editor
uses the configured font, highlighting and wrapping. Save writes that file;
Save As writes a separate copy. A local copy using the original package name
provides a document-local package override. Package-source
saves are independent of Apply in the preamble form.

Editable preamble commands are initially collapsed above the body. Expand them
to edit definitions in orange boxes. Source fragments such as preamble and
package files retain their original contents and use an editable source view.

\begin{enumerate}
\item Choose Insert, Commands, Define command. An orange box appears at the caret.
Enter the command name in the selected field.
\item Edit its body beside the assignment sign. Text and mathematics render in
place; the source toggle exposes literal code. For mixed content, use \verb|$...$|.
The minus button collapses expanded source inside the orange box.
Hover over the definition name and click its plus button to edit the complete
definition command, including its declaration type and argument syntax.
\item Use Configure definition arguments to add, remove or reorder arguments,
choose Required or Optional, and edit defaults. Right-clicking an argument marker
also changes its type. Its optional default is editable in the bracketed field.
\item Typing a new parameter such as \verb|#2| proposes a matching argument.
It is retained when editing finishes if still used. Existing declared arguments
are removed only explicitly.
\item Use Renew command to change an earlier definition, Provide to define only
when absent, or Declare math operator for an upright function name.
\end{enumerate}

\texttt{custom-commands.tex} demonstrates plain \TeX{}, optional arguments,
body redefinitions and macros imported from an adjacent package. To edit a call,
choose Home, Editing, Properties. Clicking an omitted optional
argument exposes its default; a value equal to the default is omitted from source.

\section{References, bibliography and appendices}
Insert a label at its target, then use References to choose that label and the
reference command. Equation~\eqref{eq:guide} uses \verb|\eqref|. Renaming a label
updates its references. Click an existing reference to choose a different target.
Ctrl+hover underlines a link; Ctrl+click follows it in either continuous or paged layout.
Use \verb|\pageref| for a target's native page number and \verb|\nameref| for its
heading. Page references update when the native layout changes.

In References, Bibliography, choose database files, backend and style. Installed
styles are discovered on your computer. Insert the bibliography list in the body,
then choose bibliography entries through the citation dialog; it also edits notes.
Use \texttt{references.tex} for BibTeX and \texttt{bibliography-biblatex.tex} for Biber.
Choose the preamble matching the selected bibliography backend. References, Contents, inserts a table of
contents or a list of figures/tables. These lists use short captions when supplied;
Ctrl+click an entry to visit its source, or the heading to change the list title.
Contents and bibliography blocks render natively from headings and database entries. Common numeric and author--year bibliography formats, native page
numbers, indentation and dotted leaders are supported. Custom bibliography-style
programs and package output routines remain authoritative in the compiled PDF.
The title block can reuse compiled output. Their text is selectable/copyable, but generated
entries are read only; Ctrl+click follows links or opens the block's configuration.

Choose Start appendices before the first appendix heading. Subsequent headings
use letters, and counters/references follow the document class. The dashed
Appendix divider is an editing aid. \texttt{theorems.tex} combines theorem editing,
proofs, labels and appendices, with equation and theorem numbering examples.

\section{Tables, figures and diagrams}
\subsection{Tables and images}
Insert a table and choose its dimensions. Arrow keys navigate cells. A selection
across cells becomes a rectangular cell selection. The Table tab adds, deletes,
duplicates or moves rows/columns and edits border rules around selected cells.
Use properties for column formats and longtable options.
Repeated specifications such as \verb+*{5}{c|}c+ and diagonal headers made with
\verb|\diagbox{Row}{Column}| have native previews; load the diagbox package first.

Insert a figure, choose its image, then set width/height as a number and unit.
Figure tools edit its caption, caption position, label and float placement. In an
empty figure, inserting an image fills that figure. Include PDF pages lets you
choose a PDF and page list: \verb|-| means all pages, \verb|2,4-6| selects pages,
and \verb|3-1| reverses a range. Click included pages to change their configuration.
Select a cell rectangle and choose \emph{Table, Merge selected cells};
\emph{Split cell} asks for the number of rows and columns, including for an
ordinary cell. It keeps the text in the first cell and preserves neighboring
cells with spans. Choose whether to add borders between the new cells.
Editing a row or column through a merged area splits the affected spans first.
The native renderer supports \verb|\multirow|, \verb|\multicolumn|,
array column declarations and booktabs rules/spacing. Common adjustbox width,
height, scale, frame and alignment options are supported; arbitrary package
programming still requires the compiled preview.
Native tables also support flexible \texttt{tabularx} columns, multiline
\verb|\makecell|/\verb|\thead| cells, and row, column and cell colors. Common
\texttt{enumitem} list labels, starting numbers, margins and spacing are read
from environment options or \verb|\setlist|. Saved boxes made with
\verb|\sbox| or \verb|\savebox| can be reused with \verb|\usebox|;
\verb|\typeout| stays an editing-only source object.
Try \texttt{tables-and-figures.tex} and \texttt{pdf-pages.tex}.

\subsection{TikZ and PGFPlots}
Open \texttt{tikz-plots.tex}, click a drawing and open its drawing editor.
Use the icon palettes to add objects or choose a sketch tool. Hover for each name.
Drag the divider to resize the controls; scroll the left pane to reach lower controls.
\begin{enumerate}
\item Choose Node, Path, Arrow, Rectangle or Circle and draw in the preview.
Paths accept successive clicks; Enter or double-click finishes. Backspace removes
the last point; Escape cancels. Rectangle uses opposite corners; circle uses center/radius.
\item Use Move to drag nodes, point handles or whole line segments. Choose a snap
increment when adding points. Each completed change has undo support.
\item Select an object to edit its ordered path/node rows and options, then click
Update. The up/down arrows reorder objects within their environment.
Node and path styles are separate from tool defaults.
\item Elements shows the selected object's properties directly. Edit functions and
coordinate tables there. Add, delete and reorder
datavisualization datasets; column controls appear for tables. Data set opens its name,
format and options, with a table for variables and functions. Add, remove or
reorder declarations there; the Source tab retains literal code and comments.
\item Browse options by category. Edit suboptions opens nested keys while preserving
the enclosing option. Suggestions follow the selected command, environment or suboption;
they come from the installed packages and loaded libraries. Choices and defaults appear
beside the selected key; pending values are saved when switching keys. Custom keys remain possible.
Environments edits pictures, axes and scopes;
select tikzpicture to configure scale and node defaults.
\item In Native PasTeX mode, click supported node, tick, legend or axis-label text
to edit it. Apply saves the picture; closing with unapplied changes prompts first.
\end{enumerate}
Settings enables diagram preview and selects Cached TeX or Native PasTeX. Cached
TeX uses the package engine and is the default. Native PasTeX avoids compilation
for supported objects. The \TeX{} backend handles programmable package features.
The example includes circuits, data visualization, grids, graphs, decorations,
polar plots, surfaces and colorbars for comparing both backends.

\texttt{commutative-diagrams.tex} covers TikZ-cd and XY-pic. Click supported cells
or arrow labels to edit their mathematics; expand source for other changes.
\texttt{common-packages.tex} covers units, derivatives, hyperlinks, colored boxes,
listings and longtables. Contextual Box and Listing properties group their options.

The \texttt{custom-commands.tex} example also demonstrates scoped definitions,
\verb|\begingroup|/\verb|\endgroup|, literal-name \verb|\ifcsname| and simple
character comparisons with \verb|\if|. Complex TeX conditionals remain source
objects. Native abstracts and xcolor named colors, mixtures and RGB/HTML models
use the native editor.

\section{Projects, navigation and appearance}
Open \texttt{subfiles/main.tex} to edit linked chapters directly in continuous
or page view. Save writes changes to their original files and preserves the
master's input commands. Double-dashed separators identify each subfile;
click its [hide]/[show] control to fold or unfold the linked content. Ctrl+click a separator to open its tab. These guides belong to the editor overlay. The status bar identifies the current subfile.
Turn off \emph{View, Subfiles in place} to edit the master input commands;
Ctrl+click a filename to open its tab. Relative paths resolve from the containing
file. The Methods chapter demonstrates three linked levels, including a
standalone subfile and an ordinary input. Nested files fold independently;
\verb|\include| and \verb|\subfileinclude| retain their page breaks.
Engine conditionals follow the selected compiler, including a file's
\verb|% !TeX program| directive. Switching compilers preserves unsaved subfile edits.
Empty, missing or recursive inputs remain visible commands. Make each editing operation within
one file; replacements spanning file boundaries are rejected. View changes
reset the editing history. Save checks for conflicting external changes.
Use \texttt{chinese.tex} with XeLaTeX for bilingual editing and CJK font examples.

View toggles the Navigation panel. Its heading tree and lists of figures,
tables, labels and subfiles jump to source positions in either layout. Heading
numbers follow the document counters, numbering depth and appendix settings. The zoom slider and
zoom percentage are at the right of the status bar. Click the percentage
to enter a scale or choose a fitting mode. Fit Width follows the
window size and includes nonempty margin notes in continuous layout. Fit Page
fits a paper page in paged layout. Right-click the status bar to choose page, position, word count and command
trail details. Its view icons switch editing mode and continuous or paged layout.
The fullscreen button beside the zoom controls enters or leaves focus mode;
F11 toggles it and Escape exits.
Right-click document content for its settings. Selected text offers formatting;
equations, tables, figures, references and rules offer their own editing commands.
Undo and Redo remain in the quick-access bar and Edit menu.

Settings groups appearance, source, preview and \TeX{} preferences. Choose Light,
Dark or Follow system. Set an installed source font, size, wrapping and syntax colors.
The quick-access customization dialog adds/removes commands and moves them up/down.
Document information reports path, size and word count. Help, About includes
program and detected-tool information useful for reporting a problem.

\section{When something looks wrong}
\begin{description}
\item[Compilation fails] Read the first compiler error. Check matching braces,
required packages, image/database paths and the selected engine. Source remains editable.
\item[Missing references] Allow compilation and bibliography processing
to finish. Check label spelling and that BibTeX or Biber is installed.
\item[A native drawing differs] Compare with Ctrl+R. Select Cached TeX for full
package output; keep a small source example when reporting a native-rendering issue.
\item[Fonts or symbols are missing] Check the document's font packages and installed
fonts. Chinese documents need the appropriate engine and CJK setup.
\item[Large files respond slowly] Keep source/drawing panes collapsed when unused
and avoid starting duplicate external compilations. Include your widget set,
engine, example and steps when reporting a delay.
\end{description}

\appendix
\section{Source and result examples}
The following examples pair editable source with compiled output.
Use the topic files indexed in \texttt{docs/developer-guide.tex} for larger examples.

\subsection{Text styles and literal symbols}
\begin{verbatim}
\textrm{Roman}, \textsf{sans serif}, \texttt{monospace};
\textbf{bold}, \textit{italic}, \textsc{small capitals}.
\textsf{\textit{\textbf{Nested styles}}}
{\small Small text} and {\Large large text}.
\end{verbatim}
\result
\textrm{Roman}, \textsf{sans serif}, \texttt{monospace};
\textbf{bold}, \textit{italic}, \textsc{small capitals}.
\textsf{\textit{\textbf{Nested styles}}}
{\small Small text} and {\Large large text}.

\begin{verbatim}
50\%, \$5, A\&B, \#1, file\_name, \{group\}.
\textless{}tag\textgreater{} and \textbackslash{}.
\end{verbatim}
\result 50\%, \$5, A\&B, \#1, file\_name, \{group\}.
\textless{}tag\textgreater{} and \textbackslash{}.

\subsection{Matrices, delimiters and operators}
\begin{verbatim}
\begin{pmatrix}1&2\\3&4\end{pmatrix},\quad
\left\lVert\frac{x}{y}\right\rVert,\quad
\sum_{n=1}^{\infty}\frac{1}{n^2},\quad \int_0^1 x\,dx
\end{verbatim}
\result
\[\begin{pmatrix}1&2\\3&4\end{pmatrix},\quad
\left\lVert\frac{x}{y}\right\rVert,\quad
\sum_{n=1}^{\infty}\frac{1}{n^2},\quad\int_0^1x\,dx\]

\subsection{Command defaults}
\begin{verbatim}
\def\myparen#1{\left(#1\right)}
\newcommand{\hello}[2][Hello]{#1, \textbf{#2}!}
$\myparen{\frac{1}{2}}$; \hello{reader}; \hello[Welcome]{friend}
\end{verbatim}
\def\myparen#1{\left(#1\right)}
\newcommand{\hello}[2][Hello]{#1, \textbf{#2}!}
\result $\myparen{\frac{1}{2}}$; \hello{reader}; \hello[Welcome]{friend}

\subsection{Theorems and tables}
\begin{verbatim}
% Preamble: \usepackage{amsthm}
%           \newtheorem{theorem}{Theorem}[section]
\begin{theorem}[A sum]\label{thm:sum}
For $n\geq1$, $1+\cdots+n=n(n+1)/2$.
\end{theorem}
\begin{proof}Induction on $n$.\end{proof}
\end{verbatim}
\result
\begin{theorem}[A sum]\label{thm:sum}
For $n\geq1$, $1+\cdots+n=n(n+1)/2$.
\end{theorem}
\begin{proof}Induction on $n$.\end{proof}

\begin{verbatim}
\begin{tabular}{lc}
\toprule Name & Value \\\midrule
Alpha & 1 \\ Beta & 2 \\\bottomrule
\end{tabular}
\end{verbatim}
\result
\begin{center}\begin{tabular}{lc}
\toprule Name&Value\\\midrule Alpha&1\\Beta&2\\\bottomrule
\end{tabular}\end{center}

\subsection{Shortcut reference}
\begin{center}\begin{tabular}{ll}
\toprule Shortcut&Action\\\midrule
Ctrl+O / Ctrl+S&Open / save\\
Ctrl+Shift+S / Ctrl+W&Save as / close tab\\
Ctrl+Z / Ctrl+Y&Undo / redo\\
Ctrl+X / Ctrl+C / Ctrl+V&Cut / copy / paste\\
Alt+=&Toggle inline mathematics\\
Ctrl+L&Toggle selected source in place\\
Ctrl+Shift+U&Full document source\\
Ctrl+R&Open compiled PDF\\
Ctrl+/&Comment / uncomment\\
Ctrl+Shift+N&Toggle navigation panel\\
Ctrl+F1&Collapse / restore ribbon\\
Alt+Q&Find a command\\
\bottomrule\end{tabular}\end{center}
\end{document}
